% Part: proof-theory
% Chapter: cut-elimination
% Section: ce-topmost

\documentclass[../../../include/open-logic-section]{subfiles}

\begin{document}

\olfileid[tr]{pt}{cut}{top}

\olsection{En Üst Kesmelerin Kaldırılması}

\begin{defn}
Şu !!a{proof}~$\pi$'yi ele alın; bir \CutCS{} kuralıyla biter:
\[
\AxiomC{}
\RightLabel{$\pi_1$}
\Deduce$\Gamma \fCenter \Delta, !A$
\AxiomC{}
\RightLabel{$\pi_2$}
\Deduce$!A, \Gamma \fCenter \Delta$
\RightLabel{\CutCS}
\BinaryInf$\Gamma \fCenter \Delta$
\DisplayProof
\]
ve bunun başka hiçbir \CutCS{} çıkarımı içermediğini varsayın.
$\pi$'nin \emph{kesme yüksekliği}, $\cheight{\pi} =
\pheight{\pi_1} + \pheight{\pi_2}$'dir. $\pi$'nin \emph{kesme
derecesi}, $\cutrank{\pi} = \depth{!A}$'dır.
\end{defn}

\begin{lem}\ollabel{lem:cut-adm-G3c}
\CutCS{} kuralı \Log{G3c}'de kabul edilebilirdir. Başka bir deyişle,
$\pi$ tek bir~\CutCS{} ile biten ve bunun dışında yalnızca
$\Log{G3c}$ kurallarını kullanan !!a{proof} ise aynı son-sekantın \Log{G3c}'de
!!a{proof}{}ı~$\pi'$ vardır.
\end{lem}

\begin{proof}
İspat, kesmenin !!{height}~değeri ile kesme derecesi üzerinde çift tümevarımladır.
!!{proof}~$\pi$ şöyle biter:
\[
\AxiomC{}
\RightLabel{$\pi_1$}
\Deduce$\Gamma \fCenter \Delta, !A$
\AxiomC{}
\RightLabel{$\pi_2$}
\Deduce$!A, \Gamma \fCenter \Delta$
\RightLabel{\CutCS}
\BinaryInf$\Gamma \fCenter \Delta$
\DisplayProof
\]

Tümevarım tabanı: $\cheight{\pi} = 0$ ve $\cutrank{\pi} = 0$.
$\cheight{\pi} = \pheight{\pi_1} + \pheight{\pi_2} = 0$ olduğundan
$\Gamma \Sequent \Delta, !A$ ile $!A, \Gamma \Sequent \Delta$'nın
ikisi de aksiyomdur. İki durum vardır: (a)~$!A$ bunlardan birinde
asıl değildir ya da (b)~$!A$ ikisinde de asıldır. Aşağıda her iki
durumda da $\Gamma \Sequent \Delta$'nın bir aksiyom olmak zorunda
olduğunu göstereceğiz (bunun için tümevarım varsayımına ihtiyacımız
yoktur).

Şimdi $\pi_1$ ile $\pi_2$'nin aksiyom olup olmamasına ya da bir
çıkarımla bitip bitmemesine ve $!A$'nın bu çıkarımda asıl olup
olmamasına göre bir dizi olası durumu ayırt edeceğiz. A ve~B
durumları tümevarım tabanını kapsar. Tümevarım adımında
$\cheight{\pi} > 0$ ya da $\cutrank{\pi} > 0$ (veya her ikisi)
olduğunu varsayarız. Tümevarım varsayımını kullanabiliriz: Tek
bir~\CutCS{} ile biten herhangi bir !!{proof} içindeki son~\CutCS{},
ya kesme derecesi $= \cutrank{\pi}$ ve kesmenin !!{height}~değeri
$< \cheight{\pi}$ ise ya da kesme derecesi~$< \cutrank{\pi}$ ise
giderilebilir. $\cheight{\pi} = 0$ olduğu durum (iki öncülün de
aksiyom olması) A ve~B durumlarınca kapsanır. $\cheight{\pi} > 0$
olduğu durumlar C--D durumlarınca kapsanır.

\paragraph{Durum A:}
Bir $!B$, hem $\Gamma$'da hem de~$\Delta$'da geçer.
Bu durumda $\Gamma \Sequent \Delta$ bir aksiyomdur ($!B$~atomsal ise)
ya da \olref[seq][ex]{prop:G3c-axioms} uyarınca \CutCS{} içermeyen
bir \Log{G3c}-!!a{proof}{}ı~$\pi'$ vardır. Bu, tümevarım tabanının
(a)~durumunu kapsar: $!A$, $\Gamma \Sequent \Delta, !A$ içinde ya da
$!A, \Gamma \Sequent \Delta$ içinde asıl değilse ve bunlar
aksiyomsa başka bir atomsal~$!B$ hem $\Gamma$'da hem de~$\Delta$'da
geçmek zorundadır. Öncüllerden biri $!A$'nın asıl olmadığı bir
aksiyomsa (öteki öncül bir kuralın sonucu olsa bile) bu, tümevarım
adımını da kapsar.

\paragraph{Durum B:} İki öncül de aksiyomdur; $!A$ asıldır ve
dolayısıyla atomsaldır. Sol öncül $\Gamma \Sequent \Delta, !A$'yı
ele alın. $!A$ asıl olduğundan $\Gamma$'da geçmek zorundadır (aksi
hâlde $\Gamma \Sequent \Delta, !A$ bir aksiyom olmazdı). Benzer
biçimde $!A$, $\Delta$'da geçmek zorundadır; yoksa
$!A, \Gamma \Sequent \Delta$ bir aksiyom olmazdı. Öyleyse $!A$
atomsaldır ve hem $\Gamma$'da hem de $\Delta$'da geçer; yani
$\Gamma \Sequent \Delta$ bir aksiyomdur. Bu, tümevarım tabanının
(b)~durumunu kapsar.

\begin{center}
Alternatif A ve B durumları
\end{center}

\paragraph{Durum A:} Kesmenin öncüllerinden biri
$!C, \Gamma' \Sequent \Delta', !C$ biçiminde bir aksiyomdur ve
$!C$~atomsaldır.
$!C$, kesme !!{formula}{}ü~$!A$ değilse $!C$ hem $\Gamma$'da hem
de~$\Delta$'da geçmek zorundadır. Bu durumda
$\Gamma \Sequent \Delta$ bir aksiyomdur ve bunu~$\pi'$ olarak
alabiliriz. $!C \ident !A$ kesme !!{formula}{}üyse, ya (sol öncül
aksiyomsa) $!A \in \Gamma$ ve $\Gamma = \Gamma', !A$ ya da (sağ
öncül aksiyomsa) $!A \in \Delta$ ve $\Delta = \Delta', !A$ olur.
İlk durumda $\pi_2$, $!A, !A, \Gamma' \Sequent \Delta$'nın
!!a{proof}{}ıdır ve \LeftR{\Contraction}'ın kabul edilebilirliği
(\olref[seq][inv]{prop:G3c-cont-adm}) sayesinde
$\Gamma \Sequent \Delta$'nın !!a{proof}{}ı vardır. İkinci durumda
$\pi_1$, $\Gamma \Sequent \Delta', !A, !A$'nın !!a{proof}{}ıdır ve
\RightR{\Contraction}'ın kabul edilebilirliği sayesinde
$\Gamma \Sequent \Delta$'nın !!a{proof}{}ını elde ederiz.

\paragraph{Durum B:} Öncüllerden biri
$\lfalse, \Gamma' \Sequent \Delta'$ biçiminde bir aksiyomdur.
Sol öncül böyle bir aksiyomsa $\lfalse \in \Gamma$ olur ve
$\Gamma \Sequent \Delta$ da bir aksiyomdur. Sağ öncül böyle bir
aksiyomsa ya $\lfalse \in \Gamma$ olur (bu durumda
$\Gamma \Sequent \Delta$ yine bir aksiyomdur) ya da
$!A \ident \lfalse$ olur. İkinci durumda $\pi_1$,
$\Gamma \Sequent \Delta, \lfalse$'ın !!a{proof}{}ıdır.
Bunu~$\Gamma \Sequent \Delta$'nın !!a{proof}{}ına, bir $\lfalse$
kopyasını $\pi_1$'deki her sekantın ardılından silerek dönüştürebiliriz.
(Alıştırma: Bunu~$\pheight{\pi_1}$ üzerinde tümevarımla doğrulayın.)

Şimdi ne A ne de B durumunun geçerli olduğunu varsayın. Geri kalan
durumlarda $\cheight{\pi} > 0$'dır; yani öncüllerden en az biri bir
çıkarımın sonucudur.

\paragraph{C ve D Durumları: Kesmelerin yerini değiştirme.}
C ve~D durumlarında ele alınan olasılıkların hepsinde ortak bir örüntü
vardır. \CutCS{} ile biten bir !!a{proof}{}la başlarız; bunun bir öncülü
bir~$R$ kuralıyla elde edilir ve bu kuralda kesme
!!{formula}{}ü~$!A$ asıl değildir (başka bir !!{formula}~$!D$ asıldır). Bu
!!{proof} içinde \CutCS{}, $R$ kuralını izler ve şematik olarak şöyle
görünür:
\[
\AxiomC{}
\RightLabel{$\pi_1$}
\Deduce$\Gamma \fCenter \Delta', !D, !A$
\AxiomC{}
\RightLabel{$\pi_3$}
\Deduce$!A, \Gamma, \Pi \fCenter \Delta', \Lambda$
\RightLabel{$R$}
\UnaryInf$!A, \Gamma \fCenter \Delta', !D$
\RightSubproofLabel{$\pi_2$}
\RightLabel{\CutCS}
\BinaryInf$\Gamma \fCenter \Delta', !D$
\DisplayProof
\]
Bu yalnızca $R$'nin tek öncüllü bir sağ kural olduğu, $A$'nın
$R$'nin asıl !!{formula}{}ü olmadığı ve \emph{asıl} olan
!!{formula}~$!D$'nin \CutCS{}'nin sağ öncülünün ardılında geçtiği
durumu kapsar. $!D$ öncülde geçseydi (yani $R$ bir sol kural olsaydı)
ya da \CutCS{}'nin sol öncülünde geçseydi görünüş elbette farklı
olurdu. $\Pi$ ile~$\Lambda$'daki !!{formula}s, $R$ kuralının etkin
!!{formula}s{}ini temsil eder.

Bu sırayı tersine çevirip $R$'nin \CutCS{}'yi izlediği bir
!!a{proof}{}ı ele alırız:
\[
\AxiomC{}
\RightLabel{$\pi_1'$}
\Deduce$\Gamma, \Pi \fCenter \Delta', !D, \Lambda, !A$
\AxiomC{}
\RightLabel{$\pi_3'$}
\Deduce$!A, \Gamma, \Pi \fCenter \Delta', !D, \Lambda$
\RightLabel{\CutCS}
\BinaryInf$\Gamma, \Pi \fCenter \Delta', !D, \Lambda$
\RightLabel{$R$}
\UnaryInf$\Gamma \fCenter \Delta', !D, !D$
\DisplayProof
\]
$\pi_1'$ ile $\pi_3'$yü sırasıyla $\pi_1$ ile $\pi_3$'ten elde etmek
için !!{height}-koruyan zayıflatmayı kullanırız.

\CutCS{} ile~$R$'nin sırasını değiştirdiğimiz için buna
“bir kesmeyi $R$ üzerinden yukarı geçirmek” denir. Bunu yapmak,
\CutCS{}'nin sağ öncülüyle biten !!{proof}{}ın yüksekliğini, dolayısıyla
kesmenin !!{height}~değerini azaltır. Başka bir deyişle yeni \CutCS{} ile
biten alt-!!{proof}s{}ın !!{height}s{}inin $\pi_1$ ile $\pi_3$'ün
yüksekliklerinden büyük olmadığı varsayılabilir; böylece kesme
!!{height} azalır (sağdan bir çıkarımı kaldırmış oluruz). Artık
tümevarım varsayımı uygulanır ve \CutCS{} çıkarımının
giderilebileceğini söyler. Son olarak fazla~$!D$ kopyasını
$\RightR{\Contraction}$'ın kabul edilebilirliğiyle gideririz.

\paragraph{Durum C:}
$\pi_1$ ile $\pi_2$'den en az biri tek öncüllü bir~$R$ kuralıyla
biter ve $!A$ bu kuralda asıl değildir. $!A$'nın \CutCS{}'nin sol mu sağ mı
öncülünde asıl olmadığına ve dokuz tek öncüllü kuraldan
(\LeftR{\lnot}, \RightR{\lnot}, \RightR{\lor}, \LeftR{\land},
\RightR{\lif} ve dört niceleyici kuralı) hangisinin $R$ olduğuna göre
toplam 18 durum vardır. Yalnızca iki örneği ele alacağız: $!A$'nın,
$\pi_2$'nin $\RightR{\lif}$ ya da \LeftR{\lexists} ile biten son
çıkarımında asıl olmadığı durumlar. Öteki durumlar benzerdir ve
alıştırma olarak bırakılmıştır.

$\pi$'nin şöyle bittiğini varsayın:
\[
\AxiomC{}
\RightLabel{$\pi_1$}
\Deduce$\Gamma \fCenter \Delta', !B \lif !C, !A$
\AxiomC{}
\RightLabel{$\pi_3$}
\Deduce$!A, !B, \Gamma \fCenter \Delta', !C$
\RightLabel{$\RightR{\lif}$}
\UnaryInf$!A, \Gamma \fCenter \Delta', !B \lif !C$
\RightSubproofLabel{$\pi_2$}
\RightLabel{\CutCS}
\BinaryInf$\Gamma \fCenter \Delta', !B \lif !C$
\DisplayProof
\]
burada $\Delta = \Delta', !B \lif !C$'dir. $\pi_1$'e
!!{height}-koruyan zayıflatmayı (\olref[seq][adm]{prop:weak-G3c-adm})
uygulayıp !!a{proof}~$\pi_1'$nü
$!B, \Gamma \fCenter \Delta', !B \lif !C, !C, !A$'nın ispatı olarak
elde ederiz (solda $!B$, sağda $!C$ ile zayıflatırız). Benzer biçimde
\olref[seq][adm]{prop:weak-G3c-adm} önermesini~$\pi_3$'e uygulayıp
!!a{proof}~$\pi_3'$nü $!A, !B, \Gamma \fCenter
\Delta', !B \lif !C, !C$'nin ispatı olarak elde ederiz
(sağda $!B \lif !C$ ile zayıflatırız). Tümevarım varsayımı,
\[
\AxiomC{}
\RightLabel{$\pi_1'$}
\Deduce$!B, \Gamma \fCenter \Delta', !B \lif !C, !C, !A$
\AxiomC{}
\RightLabel{$\pi_3'$}
\Deduce$!A, !B, \Gamma \fCenter \Delta', !B \lif !C, !C$
\RightLabel{\CutCS}
\BinaryInf$!B, \Gamma \fCenter \Delta', !B \lif !C, !C$
\DisplayProof
\]
kesme !!{formula}{}ü~$!A$ olmak üzere uygulanır: Bu
!!{proof}, $\pi$ ile aynı kesme derecesine ve daha düşük kesme
!!{height}~değerine sahiptir. Böylece !!a{proof}~$\pi_4$, \Log{G3c} içinde
\CutCS{} içermeksizin $!B, \Gamma \fCenter
\Delta', !B \lif !C, !C$'yi ispatlar. $\RightR{\lif}$ uygulayıp
$\Gamma \fCenter \Delta', !B \lif !C, !B \lif !C$'nin !!a{proof}{}ını
elde ederiz:
\[
\AxiomC{}
\RightLabel{$\pi_4$}
\Deduce$!B, \Gamma \fCenter \Delta', !B \lif !C, !C$
\RightLabel{$\RightR{\lif}$}
\UnaryInf$\Gamma \fCenter \Delta', !B \lif !C, !B \lif !C$
\DisplayProof
\]
!!a{proof}~$\pi'$nü $\Gamma \Sequent \Delta$'nın, yani
$\Gamma \fCenter \Delta', !B \lif !C$'nin ispatı olarak elde etmek
için \olref[seq][inv]{prop:G3c-cont-adm} önermesini
(büzülmenin kabul edilebilirliğini) uygulayın.

Şimdi $\pi$'nin şöyle bittiğini varsayın:
\[
\AxiomC{}
\RightLabel{$\pi_1$}
\Deduce$\lexists[x][!B(x)], \Gamma' \fCenter \Delta, !A$
\AxiomC{}
\RightLabel{$\pi_3$}
\Deduce$!A, !B(c), \Gamma' \fCenter \Delta$
\RightLabel{$\LeftR{\lexists}$}
\UnaryInf$!A, \lexists[x][!B(x)], \Gamma' \fCenter \Delta$
\RightSubproofLabel{$\pi_2$}
\RightLabel{\CutCS}
\BinaryInf$\lexists[x][!B(x)], \Gamma' \fCenter \Delta$
\DisplayProof
\]
Tümevarım varsayımını yine
\[
\AxiomC{}
\RightLabel{$\pi_1[!B(c) \Sequent {}]$}
\Deduce$\lexists[x][!B(x)], !B(c), \Gamma' \fCenter \Delta, !A$
\AxiomC{}
\LeftLabel{$\pi_3[\lexists[x][!B(x)] \Sequent {}]$}
\Deduce$!A, \lexists[x][!B(x)], !B(c), \Gamma' \fCenter \Delta$
\RightLabel{\CutCS}
\BinaryInf$\lexists[x][!B(x)], !B(c), \Gamma' \fCenter \Delta$
\RightLabel{\LeftR{\lexists}}
\UnaryInf$\lexists[x][!B(x)], \lexists[x][!B(x)], \Gamma' \fCenter \Delta$
\DisplayProof
\]
yukarıdaki türetime uygularız. Özdeğişken koşulunun sağlandığına dikkat edin:
$\pi_2$ içindeki~\LeftR{\lexists}'in özdeğişken koşulu, $c$'nin
$\lexists[x][!B(x)]$, $\Gamma'$ ya da~$\Delta$ içinde geçmediğini
güvence altına alır. Son olarak
\olref[seq][inv]{prop:G3c-cont-adm} önermesini uygulayın.


\paragraph{Durum D:}
$\pi_1$ ile $\pi_2$'den en az biri iki öncüllü bir~$R$ kuralıyla
biter ve $!A$ bu kuralda asıl değildir. Altı durum vardır. $!A$'nın,
$\pi_2$'nin $\RightR{\land}$ ile biten son çıkarımında asıl olmadığı
durumu ele alırız; öteki durumlar benzerdir.

$\pi$'nin şöyle bittiğini varsayın:
\[
\AxiomC{}
\RightLabel{$\pi_1$}
\Deduce$\Gamma \fCenter \Delta', !B \land !C, !A$
\AxiomC{}
\RightLabel{$\pi_3$}
\Deduce$!A, \Gamma \fCenter \Delta', !B$
\AxiomC{}
\RightLabel{$\pi_4$}
\Deduce$!A, \Gamma \fCenter \Delta', !C$
\RightLabel{$\RightR{\land}$}
\BinaryInf$!A, \Gamma \fCenter \Delta', !B \land !C$
\RightSubproofLabel{$\pi_2$}
\RightLabel{\CutCS}
\BinaryInf$\Gamma \fCenter \Delta$
\DisplayProof
\]
Tümevarım varsayımı şu !!{proof}s için uygulanır:
\[
\AxiomC{}
\RightLabel{$\pi_1'$}
\Deduce$\Gamma \fCenter \Delta', !B \land !C, !B, !A$
\AxiomC{}
\RightLabel{$\pi_3'$}
\Deduce$!A, \Gamma \fCenter \Delta', !B \land !C, !B$
\RightLabel{\CutCS}
\BinaryInf$\Gamma \fCenter \Delta', !B \land !C, !B$
\DisplayProof
\]
ve
\[
\AxiomC{}
\RightLabel{$\pi_1''$}
\Deduce$\Gamma \fCenter \Delta', !B \land !C, !C, !A$
\AxiomC{}
\RightLabel{$\pi_4'$}
\Deduce$!A, \Gamma \fCenter \Delta', !B \land !C, !C$
\RightLabel{\CutCS}
\BinaryInf$\Gamma \fCenter \Delta', !B \land !C, !C$
\DisplayProof
\]
Buradaki \Log{G3c}-!!{proof}s~$\pi_1'$ ile $\pi_1''$, $\pi_1$'den
!!{height}-koruyan zayıflatmayla
(\olref[seq][adm]{prop:weak-G3c-adm}) elde edilir (bir kez sağda
$!B$, bir kez de $!C$ ile zayıflatılır). !!{proof}s~$\pi_3'$ ile
$\pi_4'$ de $\pi_3$ ile~$\pi_4$'ten, sırasıyla $\pi_3$ ile~$\pi_4$'e
!!{height}-koruyan zayıflatma uygulanarak (sağda $!B \land !C$ ile
zayıflatılarak) elde edilir.

Tümevarım varsayımı bize \Log{G3c}-!!{proof}s{}ı~$\pi_5$ ve~$\pi_6$'yı,
sırasıyla $\Gamma \fCenter \Delta', !B \land !C, !B$ ile
$\Gamma \fCenter \Delta', !B \land !C, !C$'nin ispatları olarak
verir. Bunları, $\Gamma \Sequent \Delta', !B \land !C, !B \land !C$'nin
!!a{proof}{}ını elde etmek üzere $\RightR{\land}$ kuralıyla birleştiririz:
\[
\AxiomC{}
\RightLabel{$\pi_5$}
\Deduce$\Gamma \fCenter \Delta', !B \land !C, !B$
\AxiomC{}
\RightLabel{$\pi_6$}
\Deduce$\Gamma \fCenter \Delta', !B \land !C, !C$
\RightLabel{\RightR{\land}}
\BinaryInf$\Gamma \fCenter \Delta', !B \land !C, !B \land !C$
\DisplayProof
\]
$\Gamma \Sequent \Delta$'nın, yani
$\Gamma \fCenter \Delta', !B \land !C$'nin !!a{proof}{}ını elde etmek
için \olref[seq][inv]{prop:G3c-cont-adm} önermesini (büzülmenin kabul
edilebilirliğini) uygulayın.

\paragraph{Durum E: Kesmeleri indirgeme.}
Son durum, hem $\pi_1$ hem de $\pi_2$'nin $!A$'nın asıl olduğu
çıkarımlarla bittiği koşulla ilgilidir. $!A$'nın olası her
Her !!{main operator} için bir tane olmak üzere altı durum vardır.

Her durumda kesme !!{formula}{}ü~$!A$ olan bir \CutCS{}'yi,
$!A$'nın alt-!!{formula}{}ü üzerinde, yani $!A$'nın asıl
!!{formula} olduğu çıkarımların etkin !!{formula}{}ü üzerinde bir ya
da daha çok kesmeye dönüştürürüz. Bu durumlara bu yüzden çoğu kez
kesme çıkarımının “indirgenmesi” ya da “dönüştürülmesi” denir.

\begin{enumerate}
\item $!A \ident \lnot !B$ ise $\pi$ şöyle biter:
\[
\AxiomC{}
\RightLabel{$\pi_1'$}
\Deduce$!B, \Gamma \fCenter \Delta$
\RightLabel{\RightR{\lnot}}
\UnaryInf$\Gamma \fCenter \Delta, \lnot !B$
\LeftSubproofLabel{$\pi_1$}
\AxiomC{}
\RightLabel{$\pi_2'$}
\Deduce$\Gamma \fCenter \Delta, !B$
\RightLabel{\LeftR{\lnot}}
\UnaryInf$\lnot !B, \Gamma \fCenter \Delta$
\RightSubproofLabel{$\pi_2$}
\RightLabel{\CutCS}
\BinaryInf$\Gamma \fCenter \Delta$
\DisplayProof
\]
Tümevarım varsayımı uyarınca \CutCS{}, aşağıdaki türetimden giderilebilir:
\[
\AxiomC{}
\RightLabel{$\pi_2'$}
\Deduce$\Gamma \fCenter \Delta, !B$
\AxiomC{}
\RightLabel{$\pi_1'$}
\Deduce$!B, \Gamma \fCenter \Delta$
\RightLabel{\CutCS}
\BinaryInf$\Gamma \fCenter \Delta$
\DisplayProof
\]

\item
$!A \ident (!B \lor !C)$ ise $\pi$ şöyle biter:
\[
\AxiomC{}
\RightLabel{$\pi_1'$}
\Deduce$\Gamma \fCenter \Delta, !B, !C$
\RightLabel{\RightR{\lor}}
\UnaryInf$\Gamma \fCenter \Delta, !B \lor !C$
\LeftSubproofLabel{$\pi_1$}
\AxiomC{}
\RightLabel{$\pi_2'$}
\Deduce$!B, \Gamma \fCenter \Delta$
\AxiomC{}
\RightLabel{$\pi_2''$}
\Deduce$!C, \Gamma \fCenter \Delta$
\RightLabel{\LeftR{\lor}}
\BinaryInf$!B \lor !C, \Gamma \fCenter \Delta$
\RightSubproofLabel{$\pi_2$}
\RightLabel{\CutCS}
\BinaryInf$\Gamma \fCenter \Delta$
\DisplayProof
\]
\CutCS{} ile biten şu !!{proof}{}ı ele alın:
\[
\AxiomC{}
\RightLabel{$\pi_1'$}
\Deduce$\Gamma \fCenter \Delta, !B, !C$
\AxiomC{}
\RightLabel{$\pi_2'''$}
\Deduce$!C, \Gamma \fCenter \Delta, !B$
\RightLabel{\CutCS}
\BinaryInf$\Gamma \fCenter \Delta, !B$
\DisplayProof
\]
burada $\pi_2'''$, $\pi_2''$den !!{height}-koruyan zayıflatmayla
elde edilir. Bunun kesme derecesi $< \cutr{\pi}$'dir; çünkü
$\depth{!C} < \depth{!B \lor !C}$ olur. Dolayısıyla tümevarım
varsayımı \Log{G3c} içinde !!a{proof}~$\pi_1''$nü
$\Gamma \fCenter \Delta, !B$'nin ispatı olarak verir.
\CutCS{} ile biten şu !!{proof}{}ın
\[
\AxiomC{}
\RightLabel{$\pi_1''$}
\Deduce$\Gamma \fCenter \Delta, !B$
\AxiomC{}
\RightLabel{$\pi_2'$}
\Deduce$!B, \Gamma \fCenter \Delta$
\RightLabel{\CutCS}
\BinaryInf$\Gamma \fCenter \Delta$
\DisplayProof
\]
kesme derecesi $= \depth{!B} < \depth{!B \lor !C}$'dir; böylece
tümevarım varsayımı bir kez daha uygulanır ve !!a{proof}~$\pi'$,
$\Gamma \fCenter \Delta$'nın ispatı olarak elde edilir.

\item $!A \ident (!B \land !C)$. Alıştırma.

\item $!A \ident (!B \lif !C)$ ise $\pi$ şöyle biter:
\[
\AxiomC{}
\RightLabel{$\pi_1'$}
\Deduce$!B, \Gamma \fCenter \Delta, !C$
\RightLabel{\RightR{\lif}}
\UnaryInf$\Gamma \fCenter \Delta, !B \lif !C$
\LeftSubproofLabel{$\pi_1$}
\AxiomC{}
\RightLabel{$\pi_2'$}
\Deduce$\Gamma \fCenter \Delta, !B$
\AxiomC{}
\RightLabel{$\pi_2''$}
\Deduce$!C, \Gamma \fCenter \Delta$
\RightLabel{\LeftR{\lif}}
\BinaryInf$!B \lif !C, \Gamma \fCenter \Delta$
\RightSubproofLabel{$\pi_2$}
\RightLabel{\CutCS}
\BinaryInf$\Gamma \fCenter \Delta$
\DisplayProof
\]
\CutCS{} ile biten şu !!{proof}{}ı ele alın:
\[
\AxiomC{}
\RightLabel{$\pi_1'$}
\Deduce$!B, \Gamma \fCenter \Delta, !C$
\AxiomC{}
\RightLabel{$\pi_2''$}
\Deduce$!C, \Gamma \fCenter \Delta$
\RightLabel{\CutCS}
\BinaryInf$!B, \Gamma \fCenter \Delta$
\DisplayProof
\]
Bunun kesme derecesi~$\pi$'ninkinden düşüktür. Tümevarım varsayımı
!!a{proof}~$\pi_1''$nü $!B, \Gamma \fCenter \Delta$'nın ispatı olarak
verir.
Şimdi \CutCS{} ile biten şu !!{proof}{}ı ele alın:
\[
\AxiomC{}
\RightLabel{$\pi_2'$}
\Deduce$\Gamma \fCenter \Delta, !B$
\AxiomC{}
\RightLabel{$\pi_1''$}
\Deduce$!B, \Gamma \fCenter \Delta$
\RightLabel{\CutCS}
\BinaryInf$\Gamma \fCenter \Delta$
\DisplayProof
\]
Bunun kesme derecesi de~$\pi$'nin kesme derecesinden küçüktür;
dolayısıyla tümevarım varsayımı \Log{G3c}-!!{proof}~$\pi'$nü
$\Gamma \fCenter \Delta$'nın ispatı olarak verir.
\item $!A \ident \lforall[x][!B(x)]$ ise $\pi$ şöyle biter:
\[
\AxiomC{}
\RightLabel{$\pi_1'$}
\Deduce$\Gamma \fCenter \Delta, !B(c)$
\RightLabel{\RightR{\lforall}}
\UnaryInf$\Gamma \fCenter \Delta, \lforall[x][!B(x)]$
\LeftSubproofLabel{$\pi_1$}
\AxiomC{}
\RightLabel{$\pi_2'$}
\Deduce$!B(t), \lforall[x][!B(x)], \Gamma \fCenter \Delta$
\RightLabel{\LeftR{\lforall}}
\UnaryInf$\lforall[x][!B(x)], \Gamma \fCenter \Delta$
\RightSubproofLabel{$\pi_2$}
\RightLabel{\CutCS}
\BinaryInf$\Gamma \fCenter \Delta$
\DisplayProof
\]
Tümevarım varsayımı uyarınca \CutCS{}, aşağıdaki türetimden giderilebilir:
\[
\AxiomC{}
\RightLabel{$\pi_1''$}
\Deduce$!B(t), \Gamma \fCenter \Delta, \lforall[x][!B(x)]$
\AxiomC{}
\RightLabel{$\pi_2'$}
\Deduce$!B(t), \lforall[x][!B(x)], \Gamma \fCenter \Delta$
\RightLabel{\CutCS}
\BinaryInf$!B(t), \Gamma \fCenter \Delta$
\DisplayProof
\]
burada $\pi_1''$, $\pi_1$'den ($\pi_1'$den değil) !!{height}-koruyan
zayıflatmayla elde edilir. (!!{proof}, aynı kesme derecesine ama daha
düşük kesme !!{height}~değerine sahiptir.) Böylece
!!a{proof}~$\pi_2''$, $!B(t), \Gamma \fCenter \Delta$'nın ispatı
olarak elde edilir.

!!{proof}~$\pi_1'$nün son-sekantı
$\Gamma \Sequent \Delta, !B(c)$'dir; burada $c$, bir
$\RightR{\lforall}$ çıkarımının özdeğişkenidir. Dolayısıyla $c$,
$!B(x)$, $\Gamma$ ya da~$\Delta$ içinde geçmez. !!{proof}~$\pi$'nin
düzenli olduğunu varsayarız; öyleyse $c$, yalnızca aşağıdaki
$\RightR{\lforall}$ çıkarımının özdeğişkenidir ve yalnızca onun üstünde,
yani yalnızca~$\pi_1'$ içinde geçer; ayrıca $\pi_1'$ içindeki bir
çıkarımın özdeğişkeni değildir. $c$, $\pi_2$ içinde geçmediği için
$t$ içinde de geçemez. \olref[seq][qua]{lem:subst-regular} uyarınca
!!{proof}~$\Subst{\pi_1'}{t}{c}$,
$\Gamma \Sequent \Delta, !B(t)$'nin !!a{proof}{}ıdır. Şimdi tümevarım
varsayımını şu !!{proof}{}a uygularız:
\[
\AxiomC{}
\RightLabel{$\Subst{\pi_1'}{t}{c}$}
\Deduce$\Gamma \fCenter \Delta, !B(t)$
\AxiomC{}
\RightLabel{$\pi_2''$}
\Deduce$!B(t), \Gamma \fCenter \Delta$
\RightLabel{\CutCS}
\BinaryInf$\Gamma \fCenter \Delta$
\DisplayProof
\]
(Kesme !!{formula}{}ü~$!B(t)$ olduğu için tümevarım varsayımı uygulanır;
yani !!{proof}{}ın kesme derecesi~$\pi$'ninkinden düşüktür.)
\item $!A \ident \lexists[x][!B(x)]$ durumunu alıştırma olarak
bırakıyoruz.
\end{enumerate}
\end{proof}

\begin{prob}
\olref{lem:cut-adm-G3c} ispatındaki D durumunun, $!A$'nın
$\pi_1$'in $\LeftR{\lif}$ ile biten son çıkarımında asıl olmadığı alt
durumunu doğrulayın. (Not: Bu alıştırmanın bir bölümü, neyi
ispatlamanız gerektiğini, yani bu durumda $\pi$'nin hangi biçimde
olduğunu açıkça belirtmektir.)
\end{prob}

\begin{prob}
\olref{lem:cut-adm-G3c} ispatındaki D durumunun, $!A$'nın
$\pi_1$'in $\LeftR{\lforall}$ ile biten son çıkarımında asıl olmadığı
alt durumu doğrulayın.
\end{prob}

\begin{prob}
\olref{lem:cut-adm-G3c} ispatındaki E durumunun, $!A$'nın hem
$\pi_1$'in hem de $\pi_2$'nin son çıkarımlarında asıl olduğu ve
$!A \ident (!B \land !C)$ olduğu alt durumu doğrulayın.
\end{prob}

\begin{prob}
\olref{lem:cut-adm-G3c} ispatındaki E durumunun, $!A$'nın hem
$\pi_1$'in hem de $\pi_2$'nin son çıkarımlarında asıl olduğu ve
$!A \ident \lexists[x][!B(x)]$ olduğu alt durumu doğrulayın.
\end{prob}

E durumunda kesme !!{height}~değerinin, tümevarım varsayımını uyguladığımız
ve \CutCS{} ile biten !!{proof} için özgün ispatın kesme
!!{height}~değerinden daha büyük olabileceğine dikkat edin. Örneğin
$!A \ident (!B \lif !C)$ durumunda $\pi$'yi iki \CutCS{} çıkarımı
içeren bir !!a{proof}{}a dönüştürdük:
\[
\AxiomC{}
\RightLabel{$\pi_2'$}
\Deduce$\Gamma \fCenter \Delta, !B$
\AxiomC{}
\RightLabel{$\pi_1'$}
\Deduce$!B, \Gamma \fCenter \Delta, !C$
\AxiomC{}
\RightLabel{$\pi_2''$}
\Deduce$!C, \Gamma \fCenter \Delta$
\RightLabel{\CutCS}
\BinaryInf$!B, \Gamma \fCenter \Delta$
\RightSubproofLabel{$\pi_1''$}
\RightLabel{\CutCS}
\BinaryInf$\Gamma \fCenter \Delta$
\DisplayProof
\]
$\pi_1'$ ile $\pi_2''$ arasındaki üst \CutCS{}'nin hem kesme derecesi
hem de kesme !!{height}~değeri~$\pi$'ninkilerden düşüktür. Fakat tümevarım
varsayımının uygulanmasıyla elde edilen~!!{proof}~$\pi_1''$ artık
$\pi$'den daha düşük kesme !!{height}~değerine sahip değildir. Bununla
birlikte alt \CutCS{}'nin kesme !!{formula}{}ü $!B$ olduğu için kesme
\emph{derecesi}, $\pi$'nin kesme derecesinden düşüktür.

\olref{lem:cut-adm-G3c}, tek bir en üst \CutCS{} çıkarımını bir
!!a{proof}{}tan giderebileceğimizi gösterir. Herhangi bir sayıda
\CutCS{} çıkarımını bir !!a{proof}{}tan gidermek için bunu \CutCS{} ile
biten en üst alt-!!{proof}s{}a art arda uygulayabileceğimizi
gösterebiliriz.

\begin{thm}\ollabel{thm:cut-elim-G3c-topmost} Her !!{proof}~$\pi$,
$\Log{G3c} + \CutCS$ içindeyse \Log{G3c} içindeki bir !!a{proof}{}a
dönüştürülebilir.
\end{thm}

\begin{proof}
  \CutCS{} çıkarımlarının sayısı~$n$ olmak üzere, $\pi$ için bu sayı
  üzerinde tümevarımla. $n=0$ ise yapılacak bir şey yoktur: $\pi$ zaten
  \Log{G3c} içinde !!a{proof}{}tır. Aksi hâlde en üstteki bir
  \CutCS{} çıkarımıyla biten alt-!!{proof}~$\pi'$nü ele alın.
  Bu alt-ispat, son çıkarımı olarak tek bir \CutCS{} kuralı içerir.
  \olref{lem:cut-adm-G3c} önermesini uygulayarak bunu \CutCS{}
  içermeyen bir ispat~$\pi''$ne dönüştürün ve alt-!!{proof}~$\pi'$nün
  yerine~$\pi''$nü koyun. Sonuç, $n-1$ tane \CutCS{} çıkarımı içeren
  bir !!a{proof}{}tır; dolayısıyla tümevarım varsayımı uygulanır.
\end{proof}

\end{document}
